% Chapter 25: Letters, Messages and Notices: Writing for a Reason

\chapter{Letters, Messages and Notices: Writing for a Reason}

\section{Pedagogical Purpose}
Learners have already built a single well-organised paragraph around one main idea (Chapter 24) and learned to speak and write with an audience in mind (Chapter 23). This chapter takes both skills and points them at real, functional writing: the kind of letters, messages, and notices people actually write and read every day. Unlike descriptive writing, functional writing is judged mainly by whether it does its job — informs, requests, or announces clearly, in the right format, for the right reader.

\begin{learningobjectives}
The learner can:
\begin{itemize}
    \item explain the purpose and typical reader of a letter, a message, and a notice.
    \item write an informal letter using the correct parts (date, greeting, body, closing, signature).
    \item write a short, complete message containing only the essential information.
    \item write a school notice using a heading, clear body, and issuing authority.
    \item choose the most suitable format (letter, message, or notice) for a given situation.
\end{itemize}
\end{learningobjectives}

\section{Prerequisite Check}
\begin{quickcheck}
Rewrite this blunt sentence more politely, the way you learned in Chapter 23.

1. "Come to school early tomorrow."
\end{quickcheck}

\begin{rememberbox}
You already know how to choose polite, audience-appropriate language (Chapter 23) and how to organise sentences into one clear, well-ordered paragraph (Chapter 24). This chapter uses both skills for three specific jobs: writing a letter, a message, and a notice.
\end{rememberbox}

\section{Chapter Opener}
Ms. Sharma has pinned three very different pieces of writing on the classroom board.

\textbf{Item 1 (a letter):} "Dear Nani, I hope you are keeping well. I wanted to tell you about my school sports day..."

\textbf{Item 2 (a message):} "Mum — gone to Rohan's house to finish our project. Back by 6 pm. — Maya"

\textbf{Item 3 (a notice):} "NOTICE: The school will remain closed on Friday for Founder's Day. — Principal, Greenwood School"

\textbf{Rohan:} "These all look so different! One is long and friendly, one is really short, and one sounds official."

\textbf{Ms. Sharma:} "Exactly, Rohan! Each of these does a different job for a different reader. Today we'll learn when to use each one — and how to write each one well."

\section{Language Experience}
\textbf{A Letter (to a grandparent):}
\begin{tcolorbox}[colback=white, colframe=gray!20, sharp corners, boxsep=0pt, left=5pt, right=5pt, top=5pt, bottom=5pt, fontupper=\ttfamily, breakable]
12 Lake Road,
Chennai.
14 August 2026

Dear Nani,

I hope you are keeping well. I wanted to tell you that our school sports day is next Friday, and I am running in the 100-metre race! I have been practising every evening after school. I hope you can come and watch me.

I miss you very much and hope to see you soon.

With love,
Maya
\end{tcolorbox}

\textbf{A Message (left on the kitchen table):}
\begin{tcolorbox}[colback=white, colframe=gray!20, sharp corners, boxsep=0pt, left=5pt, right=5pt, top=5pt, bottom=5pt, fontupper=\ttfamily, breakable]
Mum --- gone to Rohan's house to finish our science project. Back by 6 pm. Please don't worry! --- Maya
\end{tcolorbox}

\textbf{A Notice (on the school board):}
\begin{tcolorbox}[colback=white, colframe=gray!20, sharp corners, boxsep=0pt, left=5pt, right=5pt, top=5pt, bottom=5pt, fontupper=\ttfamily, breakable]
NOTICE
Date: 14 August 2026
The school will remain closed on Friday, 21 August, on account of Founder's Day. Classes will resume as usual on Monday.
--- Principal, Greenwood School
\end{tcolorbox}

\section{Look and Notice}
\begin{thinkbox}
\begin{enumerate}
    \item Who is each piece of writing meant for — a family member, a parent at home, or the whole school?
    \item Which one is the longest, and why might it need to be longer than the others?
    \item Which one is the shortest? What information did it have to include even though it is so short?
    \item What appears at the very end of the notice that does not appear at the end of the message?
\end{enumerate}
\end{thinkbox}

\section{Discover / Explicit Teaching}
Different situations call for different kinds of writing. Each format below has its own job and its own shape:

\begin{table}[h!]
\centering
\begin{tabular}{|l|p{4cm}|p{4cm}|l|}
\hline
\textbf{Format} & \textbf{Purpose} & \textbf{Typical Reader} & \textbf{Typical Length} \\
\hline
\textbf{Letter} & To share news, feelings, or a request personally & One person you know (often family or a friend) & Several sentences to a few short paragraphs \\
\textbf{Message} & To pass on essential information quickly & Someone who needs to know something *right now* & One to three sentences \\
\textbf{Notice} & To inform a whole group of people at once & A class, a school, or the public & A few clear lines, often with a heading \\
\hline
\end{tabular}
\end{table}

\section{Grammar Word}
\begin{grammarword}{INFORMAL LETTER}
An \textbf{informal letter} is a personal letter written to someone you know well, such as a family member or friend, using a friendly, warm tone and a fixed set of parts: the sender's address and date, a greeting, a body, a closing phrase, and a signature.

\textit{Examples:}
\begin{itemize}
    \item Greeting: "Dear Nani,"
    \item Closing: "With love, Maya" or "Your loving grandson, Rohan"
\end{itemize}

\textit{Non-example:} A letter to the Principal requesting leave would use a more formal greeting ("Respected Sir/Madam") and closing ("Yours sincerely") — this chapter focuses only on the \textbf{informal} letter, written to someone you know personally.

\textit{Quick Check:} Which closing suits an informal letter to a cousin: "Yours sincerely" or "With love"?
\end{grammarword}

\section{Core Explanation}
\subsection*{The parts of an informal letter}
Address and date (top right), greeting ("Dear \_\_\_,"), body (the message itself), closing phrase ("With love," / "Your friend,"), signature (first name only). \textit{Check:} Does the letter include all five parts, in the right order, and does the greeting match the closing in warmth?

\subsection*{Writing a complete message}
"Mum — gone to Rohan's house to finish our project. Back by 6 pm. — Maya" answers *where, why,* and *when.* A message like "Mum — gone out. — Maya" is too short — it leaves the reader worried. \textit{Check:} If I imagine the reader finding only this note, do they know everything they need to know?

\subsection*{Writing a clear notice}
A heading (often "NOTICE"), the date, a short body stating exactly what, when, and where, and the issuing authority at the end (e.g., "— Principal, Greenwood School"). \textit{Check:} Could someone glance at this notice for five seconds and still know the most important fact?

\subsection*{Choosing the right format}
Telling your whole class that the library is closed tomorrow $\rightarrow$ a **notice**. Telling your best friend you are sorry you missed their birthday $\rightarrow$ a **letter**. \textit{Check:} Am I writing to one person I know personally (letter), to whoever finds this note (message), or to a whole group at once (notice)?

\section{Worked Example}
\begin{examplebox}
\textbf{Task:} Write a short informal letter from Rohan to his friend Anya, who has moved to a new city, telling her about a class trip.

\textit{Step 1:} Choose the correct format. Rohan is writing to one person he knows personally, so this should be an \textbf{informal letter}.
\textit{Step 2:} Add the address and date at the top right.
\textit{Step 3:} Choose a warm, personal greeting: "Dear Anya,"
\textit{Step 4:} Write the body — share the news clearly, in a friendly tone.
\textit{Step 5:} Choose a closing phrase that matches the friendly tone, then sign with a first name only.

\textbf{Answer:}
\begin{verbatim}
45 Hill View,
Bengaluru.
14 August 2026

Dear Anya,

I hope you are settling in well at your new school! I wanted to tell you about our class trip to the science museum last week. We saw a real rocket engine, and Ms. Sharma let us try a hands-on experiment about magnets. I wished you were there with us.

Write back soon and tell me about your new school!

Your friend,
Rohan
\end{verbatim}
\end{examplebox}

\section{Guided Practice}
\begin{activitybox}{Activity A: Identification}
Say whether each situation calls for a letter, a message, or a notice.
\begin{enumerate}
    \item Telling the whole school that the annual day has been postponed.
    \item Telling your grandmother about your new pet.
    \item Telling your brother you've gone to a friend's house and will be back by 7 pm.
\end{enumerate}
\end{activitybox}

\begin{activitybox}{Activity B: Completion}
Add the missing part to complete this letter opening.
\begin{enumerate}
    \item "\rule{5cm}{0.4pt} / Dear Uncle Vikram, / I hope this letter finds you well..."
\end{enumerate}
\end{activitybox}

\section{Independent Practice}
\begin{activitybox}{Independent Practice}
\begin{enumerate}
    \item Write a short message from you to a family member, explaining that you will be home late from an after-school activity.
    \item Write a short school notice announcing that the school book fair will be held next Tuesday in the school hall.
    \item Write the greeting and closing (only) of an informal letter to a close friend you have not seen in a long time.
    \item \textit{Reasoning:} Why do you think a notice almost always includes a date and an issuing authority, even when the notice itself is very short?
\end{enumerate}
\end{activitybox}

\section{Common Confusion}
\begin{warningbox}
Learners often write a notice in the same warm, personal tone as a letter (for example, starting a school notice with "Hi everyone!" and ending it with "Lots of love"), or write a message so briefly that it leaves out essential details like *where* or *when*.

\textit{How to check:} For a notice, ask, "Is this addressed to a group, stated clearly, and signed with a title or role rather than a personal closing?" For a message, ask, "Have I answered who, what, where, and when?"
\end{warningbox}

\section{Writing Connection}
Choose **one** of the following and write it in the correct format, using everything you have learned in this chapter:
\begin{itemize}
    \item An informal letter to a friend or family member telling them about something exciting that happened at school recently.
    \item A school notice announcing a sports day, book fair, or holiday.
    \item A message you might leave for a family member explaining where you have gone and when you will be back.
\end{itemize}

\section{Summary}
\begin{summarybox}
\begin{itemize}
    \item A \textbf{letter} shares personal news or feelings with one known reader, using address and date, greeting, body, closing, and signature.
    \item A \textbf{message} passes on essential information quickly — it should answer who, what, where, and when.
    \item A \textbf{notice} informs a whole group at once, using a heading, a clear body, and an issuing authority.
    \item Choosing the right format depends on your \textbf{purpose} and your \textbf{audience}.
\end{itemize}
\end{summarybox}

\section{Key Terms}
\begin{table}[h!]
\centering
\begin{tabular}{|l|p{9cm}|}
\hline
\textbf{Term} & \textbf{Simple Meaning} \\
\hline
Informal letter & A personal letter written to someone you know well \\
Message & A short note giving essential information quickly \\
Notice & A short, official piece of writing informing a group \\
Issuing authority & The person or role responsible for a notice (e.g., Principal) \\
Audience & The person or people meant to read a piece of writing \\
\hline
\end{tabular}
\end{table}

\begin{selfcheck}
I can:
\begin{itemize}
    \item explain the purpose and typical reader of a letter, a message, and a notice.
    \item identify the parts of an informal letter.
    \item write a short informal letter in the correct format.
    \item write a complete, concise message.
    \item write a school notice with a heading and issuing authority.
    \item choose the correct format for a given real-life situation.
\end{itemize}
\end{selfcheck}